等待问题先明确停止是否包含触发事件的那次观测，以及是在每步还是每轮末检查。对非负整数值的停止次数 $T$，尾概率直接给出
\[
\E T=\sum_{n\ge0}\Prob(T>n),
\]
等式也允许两边为无穷；完整分布则由相邻尾概率相减得到。

独立均匀比例的负对数是指数增量，可把乘法缩减转成到达时间问题。仅涉及相邻大小关系时，独立同分布连续样本的等可能排名往往比密度更方便。正增量的越界问题可汇总各次累计和的密度。

写出一、二阶矩递推前，应确认所用矩有限；处理无限随机和或乘积的极限时，则要说明交换极限与期望的依据。
\par\smallskip
